\section{Towards stellarators: a global beam--shell model}
\label{sec:stellarator}
% This section is the forward-looking part of the paper: it states the
% method and shows first results of the global model on a tokamak TF
% D-coil (where a reference exists), and the plan for stellarators.

\subsection{Why the TF approach does not transfer directly}
Stellarator coils are non-planar and carry, besides the in-plane forces of
a TF coil, large out-of-plane and locally varying net forces. There is no
wedged inner vault: each coil is held by a coil case and by an
inter-coil support structure, and the loads cannot be reduced to a
membrane hoop tension plus a centring force. The 1D/2D radial-build
treatment of section~\ref{sec:tf} is therefore not applicable as is.

\subsection{Proposed approach}
We split the problem in two levels:
\begin{enumerate}
\item \textbf{Local level (conductor and winding pack).} The cable and the
winding pack are designed exactly as in the TF solver --- same conductor
module, same jacket sizing, same grading on the local peak field --- but
\emph{without the wedge}: the winding pack cross-section is no longer
bounded by the wedge angle $2\pi/N_\mathrm{TF}$, but by a prescribed
rectangular envelope around the filament (the coil-to-coil and
coil-to-plasma clearances from the coil optimisation).
\item \textbf{Global level (structure).} The coil and its support are
modelled with beams along the coil centreline (winding pack + case,
smeared section properties from the local level) and shells for the
inter-coil structure. Line loads come from Biot--Savart on the full coil
set. The model returns section forces and moments along the coil, which
are fed back to the local level to check jacket and case stresses.
\end{enumerate}

\subsection{Beam--shell formulation}
% - Curved 3D beam (Timoshenko) along an arbitrary space curve r(s):
%   local frame from the filament (tangent + winding-pack orientation /
%   finite-build frame used by the coil optimisation code), not Frenet
%   (ill-defined where curvature -> 0).
% - Section stiffness EA, EI_n, EI_b, GJ, GA from homogenisation of the
%   WP (K_e,rad, K_e,tor already computed in the conductor module) and
%   of the case.
% - Shell elements for inter-coil structure and, for tokamak TF, the
%   outer intercoil structure; beam-shell coupling by rigid/elastic
%   links at the support locations.
% - Loads: f(s) = I t(s) x B(r(s)), B from Biot-Savart of all coils
%   (finite build / multi-filament near the coil itself).
% - Boundary conditions: gravity supports, symmetry (field periods,
%   stellarator symmetry; N_TF sectors for a tokamak).
% - Outputs: N, V, M, T along s -> stresses in WP and case via section
%   recovery -> Tresca/von Mises checks -> iterate case thickness.
\todo{equations; choice of solver (own FE in MATLAB vs external
FE code driven by MADE).}

\subsection{First application: D-shaped TF coil}
The same model applies to the D-shaped TF coil, which is planar but not
straight: it removes the Princeton-D assumption, captures the bending
introduced by a non-ideal D (e.g. shape constrained by the build), and
provides out-of-plane loads when PF fields are included. It is the
natural verification case before moving to stellarators, since
reference FE results exist for tokamak TF coils \todo{choose reference:
EU-DEMO / ITER / DTT TF FE analyses}.
\todo{Fig. 6: TF D-coil beam model, bending moment along s vs reference.}

\subsection{Roadmap for stellarator coils}
% 1. Reader for coil filaments (e.g. coils-file format used by
%    stellarator optimisation codes) + finite-build frame.
% 2. Local level: TF chain without wedge -> WP + cable on local B_peak.
% 3. Global level: beam-shell model, loads from Biot-Savart.
% 4. Benchmark on a published stellarator coil set
%    \todo{choose: e.g. W7-X non-planar coil (published FE data)}.
\todo{status of each step at submission time.}
